\section{Proof of the comparison after oracle tuning}
\label{sec:predictive-proofs}

We prove Theorem~\ref{thm:conference-predictive} using exact error moments. We
first restrict the rate parameter in \eqref{eq:conference-oracle-rates} to
$\alpha\in[0,1]$. For each signal family, we then bound all larger rates for
the method with higher risk, showing that none can reverse the ordering. Two
independent implementations,
\texttt{experiments/check\_predictive\_theorem.py} and
\texttt{experiments/check\_predictive\_review.py}, check the rational
certificates and the reductions of the moment equations. These checks run with
\texttt{make check-predictive-theorem} and \texttt{make check-predictive-review}
in the locked environment.

The proof for bounded rates has undergone a separate AI-assisted mathematical
review. For the extension to unrestricted rates, we have local algebraic checks
and a Lean formalization of the lower bounds conditional on the moment
recurrence. Independent review of that extension and human verification of the
complete theorem remain open in this working draft.

\subsection{Exact moments}

By homogeneity, we may take $\sigma^2=1$. Put
$s=(\beta_C^*)^2$, $r=(\beta_R^*)^2$, $Q=10^{-4}$, and $\tau=Mb$.
We write $C_k=\E[(\beta_{C,k}-\beta_C^*)^2]$ and
$H_k=\E[(\beta_{R,k}-\beta_R^*)^2]$. Independence gives
\begin{equation}
 R_{R,k}=C_k/8+(1-q)^2H_k/2.
 \label{eq:predictive-risk}
\end{equation}
The required feature moments are $x_C^2=1/4$, $\E x_R^2=b$,
$\E x_R^4=b(1-3b)$, and $\E x_C^2x_R^2=b/4$. Mixed moments with an odd
power of $x_C$ vanish. Expanding the squared update and using the conditional
noise assumptions therefore gives
\begin{equation}
 \binom{C_{k+1}}{H_{k+1}}
 =\begin{pmatrix}
    A&4u^2b\\ b\alpha^2/4&D_b
   \end{pmatrix}\binom{C_k}{H_k}
   +\binom{4u^2}{b\alpha^2},\qquad
 \binom{C_0}{H_0}=\binom{s}{r},
 \label{eq:predictive-moments}
\end{equation}
where $A=(1-u)^2$, $D_b=1-b\alpha(2-\alpha+3b\alpha)$, and
$u=\alpha/4$ for U or $u=b\alpha$ for S. The diagonal equations and the risk do
not involve cross error moments, so changing the signs of the coefficients
leaves the risk unchanged. Every matrix entry is nonnegative for every
$\alpha\geq0$, since the diagonal entries are expectations of squared update
coefficients.

For $0<\alpha\leq1$, the row sums are
\[
 1-u\{2-u(1+4b)\}<1,\qquad
 1-b\alpha\{2-5\alpha/4+3b\alpha\}<1.
\]
The homogeneous moments therefore contract, and the moments contributed by noise
stay bounded. Since the cross moment is bounded by $\sqrt{C_kH_k}$, all
positive rates that we use below as candidates for the method with lower risk
are mean-square stable. At $\alpha=0$, the predictor stays at initialization.

\subsection{Dominant rare signal on the bounded rate interval}

Suppose $r=1$ and $s\leq1/100$. For U, we have $A\leq D_b$ because the
common contraction factor falls below one by at least $7\alpha/16$, whereas
the rare factor falls below one by at most $2b\alpha$. Dropping nonnegative
terms in \eqref{eq:predictive-moments} and summing the noise terms gives, with
$z=D_b^k$,
\begin{equation}
 R_{R,k}^U\geq\frac12\left[(1-q)^2z+
   \left\{\frac{(1-q)^2\alpha}{2-\alpha+3b\alpha}
                +\frac{\alpha}{8-\alpha}\right\}(1-z)\right].
 \label{eq:predictive-u-lower}
\end{equation}
For the common-noise term, we use $1-A^k\geq1-D_b^k$. The bound also holds
directly at zero rate. Since the right side is affine in $z$, we can bound its
minimum over $k\leq M$ by its values at $z=1$ and $z=D_b^M$.

Set $d(\alpha)=\alpha(2-\alpha)$, $Z=e^{-\tau d(\alpha)}$, and
\[
 N(\alpha)=\frac{\alpha}{2-\alpha}+\frac{\alpha}{8-\alpha},\qquad
 F_U^-(\alpha,\tau)=\tfrac12[Z+N(\alpha)(1-Z)].
\]
Writing $d_b=d(\alpha)+3b\alpha^2$, the logarithm series yields
\[
 0\leq-\frac{\log(1-bd_b)}b-d(\alpha)
 \leq b\left[3+\frac{(1+3Q)^2}{2\{1-Q(1+3Q)\}}\right]<4b.
\]
Hence $|D_b^M-Z|\leq16q$. The initial and noise coefficients in
\eqref{eq:predictive-u-lower} differ from $1$ and $N(\alpha)$ by at most
$2q$ and $5q$, respectively, and $|1-N(\alpha)|\leq1$. It follows that
\begin{equation}
 R_{R,k}^U\geq\min\{1/2,F_U^-(\alpha,\tau)\}-11q,
 \qquad 0\leq\alpha\leq1,\quad 0\leq k\leq M.
 \label{eq:predictive-u-transfer}
\end{equation}

For S, we use $Y_k=C_k/4+H_k$ and
$\widetilde D=1-bd(\alpha)+b^2\alpha^2$. The exact recurrence gives
\[
 Y_{k+1}=\frac{\widetilde D}{4}C_k+
          (D_b+b^3\alpha^2)H_k+b\alpha^2(1+b)
 \leq\widetilde D Y_k+b\alpha^2(1+b),
\]
since $\widetilde D-(D_b+b^3\alpha^2)=b^2\alpha^2(4-b)\geq0$.
Using $R_{R,M}^S\leq Y_M/2$ and summing the geometric series gives
\[
 R_{R,M}^S\leq\frac12\left[(r+s/4)\widetilde D^M+
       \frac{\alpha(1+b)}{2-\alpha-b\alpha}(1-\widetilde D^M)\right].
\]
Define
\[
 F_S(\alpha,\tau;s,r)=\tfrac12\left[(r+s/4)Z+
                    \frac{\alpha}{2-\alpha}(1-Z)\right].
\]
The logarithm remainder and plateau difference satisfy
\[
 |\widetilde D^M-Z|\leq4b\left[1+\frac1{2(1-Q)}\right],\qquad
 0\leq\frac{\alpha(1+b)}{2-\alpha-b\alpha}
       -\frac{\alpha}{2-\alpha}\leq\frac{2b}{1-Q}.
\]
Consequently, for the rare-signal family,
\begin{equation}
 R_{R,M}^S\leq F_S(\alpha,\tau;1/100,1)+5q.
 \label{eq:predictive-s-transfer}
\end{equation}
The unrounded coefficient of $q$ is at most
$\{(401/400)4[1+1/(2(1-Q))]+2/(1-Q)\}/2<5$.

We choose the stable candidate $\alpha_S=2/(\tau+4)\in[1/4,2/5]$ and $k=M$.
We next construct a rational certificate for the bound
\begin{equation}
 \min\{1/2,F_U^-(\alpha,\tau)\}
 -F_S(2/(\tau+4),\tau;1/100,1)\geq1/100
 \quad\text{on }[0,1]\times[1,4].
 \label{eq:predictive-certificate}
\end{equation}
Both $d$ and $N$ increase on $[0,1]$. On a rectangle
$[a_0,a_1]\times[t_0,t_1]$, we put $N_0=N(a_0)$ and bound $Z$ between
$e^{-t_1d(a_1)}$ and $e^{-t_0d(a_0)}$. Since $(N+(1-N)Z)/2$ increases with
$N$, we obtain a lower bound by using $N_0$ and the lower endpoint for $Z$ if
$N_0\leq1$, or the upper endpoint otherwise. We then take the minimum with
$1/2$. The S candidate has plateau $1/(\tau+3)$ and exponent
$4\tau(\tau+3)/(\tau+4)^2$. Its risk decreases with $\tau$, so its value at
$t_0$ gives an upper bound on the rectangle.

For rational $x\in[0,4]$, put $P_{18}(x)=\sum_{j=0}^{18}x^j/j!$. Positivity
and the ratio of successive omitted terms give
\begin{equation}
 P_{18}(x)\leq e^x\leq P_{18}(x)+\frac{x^{19}/19!}{1-x/20}.
 \label{eq:predictive-exp-certificate}
\end{equation}
Taking reciprocals gives rational enclosures for $e^{-x}$. Starting from the
whole rectangle, the theorem checker subdivides each unresolved rectangle into
four equal rectangles. It accepts a leaf only when the lower margin is at least
$1/100$. The 5,035 accepted leaves have maximum depth nine and total area
exactly three, with minimum rational margin above $0.01000005$. The independent
review checker uses a different exponential enclosure and subdivision, yielding
3,925 leaves. Both constructions certify the inequality over the continuous
domain.

Combining \eqref{eq:predictive-u-transfer},
\eqref{eq:predictive-s-transfer}, and \eqref{eq:predictive-certificate},
we find that every U candidate in the bounded rate interval has risk exceeding
that of the S candidate by at least
\[
 1/100-16q\geq0.0084>0.007.
\]

\subsection{Dominant common signal on the bounded rate interval}

Suppose $s=1$ and $r\leq1/100$. For S, the common mean error is
$-(1-b\alpha)^k\beta_C^*$. If $\alpha\leq1/4$, Jensen's inequality and the
logarithm bound give
\[
 R_{R,k}^S\geq\tfrac18(1-b\alpha)^{2k}
 \geq\tfrac18\exp\{-2/(1-Q/4)\}>1/60.
\]
If $1/4\leq\alpha\leq1$, put $z=(1-b\alpha)^k$. Since
$D_b\leq1-b\alpha$ and $\alpha/(2-\alpha+3b\alpha)\geq\alpha/2$, common
squared bias and direct rare observation noise yield
\[
 R_{R,k}^S\geq z^2/8+\frac{(1-q)^2\alpha}{4}(1-z)
 \geq z^2/8+\frac3{50}(1-z)\geq\frac{33}{625}>\frac1{60}.
\]
We obtain the penultimate inequality by minimizing the quadratic over $[0,1]$.

For U, we put $h=\sqrt b\leq1/100$, choose $\alpha_U=4h$, and stop at $M$.
Then $u=h$. Write $H_{\max}=\max_{0\leq j\leq M}H_j$. Geometric summation
in the common equation of \eqref{eq:predictive-moments} gives
\[
 C_j\leq A^j+\frac{4h}{2-h}+\frac{4h^3}{2-h}H_{\max}.
\]
To obtain an upper bound, we drop the contraction in the rare coordinate and
sum this inequality. Using $M\leq4/h^2$ and
$\sum_{j\geq0}A^j=1/\{h(2-h)\}$ then gives
\[
 H_{\max}\leq r+\frac{68h^3}{2-h}+64h^2+
                       \frac{64h^5}{2-h}H_{\max}.
\]
Thus
\[
 H_{\max}\leq H_*(h):=
 \frac{1/100+68h^3/(2-h)+64h^2}{1-64h^5/(2-h)}.
\]
Using $M\geq1/h^2$ in the common bound now gives
\[
 R_{R,M}^U\leq\frac{e^{-2/h}}8+
     \frac{h}{2(2-h)}+
     \left\{\frac12+\frac{h^3}{2(2-h)}\right\}H_*(h)
 <0.01072966<\frac{11}{1000}.
\]
Each term increases with $h$ on this interval. At $h=1/100$, the checker
verifies the bound rationally using $e^{-200}\leq2^{-200}$. It verifies the
lower exponential bound above using \eqref{eq:predictive-exp-certificate}.
Every S candidate in the bounded rate interval therefore has risk exceeding
that of the U candidate by more than
$1/60-11/1000>1/200$.

\subsection{Removing the rate cap}

For the rare-signal family, $H_0=1$. Put $m_b=1/(1+3b)$. For every
$\alpha\geq1$,
\[
 D_bm_b+b\alpha^2-m_b
   =\frac{2b\alpha(\alpha-1)}{1+3b}\geq0.
\]
Dropping the nonnegative common-to-rare term in
\eqref{eq:predictive-moments} and using induction gives $H_k\geq m_b$ at
every finite stop. Hence
\[
 R_{R,k}^U(\alpha)\geq\frac{(1-q)^2}{2(1+3b)}
 \geq\frac{(1-Q)^2}{2(1+3Q)}>\frac{49}{100}.
\]
The original S candidate remains below $1/3$, since
\[
 R_{R,M}^S(\alpha_S)
 \leq F_S(2/(\tau+4),\tau;1/100,1)+5q
 \leq F_S(2/5,1;1/100,1)+5Q<\frac13.
\]
Rational exponential enclosures certify the last bound. Every added U candidate
therefore has risk exceeding that of the S candidate by more than
$49/100-1/3>7/1000$.

For the common-signal family, $C_0=1$ and $H_0\geq0$. Under S, let
$A=(1-b\alpha)^2$. The identity
$D_b-A=b\alpha^2(1-4b)\geq0$ and positivity imply
\[
 C_k\geq A^k,\qquad
 H_k\geq b\alpha^2\sum_{j=0}^{k-1}D_b^j
       \geq b\alpha^2\sum_{j=0}^{k-1}A^j.
\]
If $A\geq1$, the common error gives $R_{R,k}^S\geq1/8$. Otherwise
$0<b\alpha<2$, and geometric summation gives
\[
 R_{R,k}^S\geq\frac{A^k}{8}
       +\frac{(1-q)^2\alpha}{2(2-b\alpha)}(1-A^k).
\]
For $\alpha\geq1$, the second coefficient is at least $(1-q)^2/4>1/8$.
Thus $R_{R,k}^S\geq1/8$ for every added rate and stop, including $k=0$.
The stable U candidate remains below $11/1000$, giving a gap greater than
$1/8-11/1000>1/200$.

For each signal family, we have now bounded every candidate of the method with
higher risk against one admissible stable candidate of the other method. Taking
the infima over the enlarged searches proves both inequalities. Restricting
each search to stable rates preserves the bounds and the candidates.
Multiplying by $\sigma^2$ restores the stated gaps. The argument requires
neither stability of all added rates nor attainment of the unrestricted infima.
